\section{Alcance de la investigación}

La presente investigación tiene un alcance correlacional, debido a que busca determinar la relación existente entre el tiempo de exposición a pantallas digitales, los hábitos posturales y la aparición de síntomas físicos en estudiantes de Ingeniería de Sistemas de la Universidad Pedagógica y Tecnológica de Colombia (UPTC).

Según Hernández-Sampieri et al., los estudios correlacionales tienen como finalidad conocer la relación o grado de asociación existente entre dos o más variables dentro de un contexto determinado [5].

El estudio no pretende modificar las condiciones de los participantes ni implementar tratamientos, sino identificar patrones y asociaciones entre las variables mediante la recopilación de información obtenida a través de instrumentos de medición.

La investigación analizará:

\begin{itemize}
    \item Tiempo diario de exposición frente a dispositivos digitales.
    \item Hábitos posturales durante las actividades académicas.
    \item Frecuencia de síntomas como fatiga visual y molestias musculoesqueléticas.
\end{itemize}

Los resultados estarán delimitados a los estudiantes del programa de Ingeniería de Sistemas de la UPTC, por lo cual no representan necesariamente a toda la población universitaria.
\estado{En construcción}
\notafase{Se redacta en la Fase 3 del curso.}
